\documentclass[11pt]{article}
\usepackage[a4paper,margin=1in]{geometry}
\usepackage{fontspec}
\usepackage[boldfont=false]{xeCJK}
\setCJKmainfont{FandolSong-Regular.otf}
\usepackage{amsmath}
\usepackage{amssymb}
\usepackage{array}
\usepackage{booktabs}
\usepackage{graphicx}
\usepackage{adjustbox}
\usepackage{enumitem}
\setlist[itemize]{topsep=2pt,partopsep=0pt,parsep=0pt,itemsep=2pt,leftmargin=1.4em}
\setlist[enumerate]{topsep=2pt,partopsep=0pt,parsep=0pt,itemsep=2pt,leftmargin=1.6em}
\usepackage{parskip}
\usepackage{fvextra}
\fvset{breaklines=true,breakanywhere=true,fontsize=\small}
\setlength{\parindent}{0pt}

\begin{document}

\begin{center}
  {\LARGE\bfseries The Periodic Table: chemical periodicity}\\[0.35em]
  {\large A Level Chemistry}
\end{center}
\vspace{0.6em}

\section*{Physical properties across Period 3}

\textbf{Periodicity} 周期性 means that properties repeat in a regular pattern as you go across each \textbf{period} 周期 of the Periodic Table. Period 3 (Na to Ar) is the standard example.

\begin{center}
\adjustbox{max width=\linewidth}{%
\begin{tabular}{ll}
\toprule
\textbf{Property} & \textbf{Trend across Period 3} \\
\midrule
\textbf{atomic radius} 原子半径 & gets smaller (more nuclear charge pulls the same shell in) \\
\textbf{ionic radius} 离子半径 & positive ions are small; from $\text{P}^{3-}$ onwards the negative ions are larger \\
\textbf{melting point} 熔点 & rises to a peak at silicon, then falls sharply \\
\textbf{electrical conductivity} 导电性 & high for Na, Mg, Al; almost zero from Si onwards \\
\bottomrule
\end{tabular}}
\end{center}

The melting point and conductivity follow from the structure and bonding:

\begin{itemize}
\item \textbf{Na, Mg, Al} are \textbf{giant metallic} 金属晶体. Melting points rise (Na → Al) because each atom gives more delocalised electrons and the ions get smaller, so the bonding is stronger. They conduct well.

\item \textbf{Si} is \textbf{giant molecular} 原子晶体 (giant covalent). It has the highest melting point, because strong covalent bonds must be broken. It barely conducts.

\item \textbf{P, S, Cl, Ar} are \textbf{simple molecular} 分子晶体 (or single atoms). Their melting points are low, because only weak intermolecular forces break. They do not conduct.

\end{itemize}

\section*{Chemical properties across Period 3}

\subsection*{Reactions with oxygen, chlorine and water}

With \textbf{oxygen}:

\[
4\text{Na} + \text{O}_2 \rightarrow 2\text{Na}_2\text{O} \qquad 2\text{Mg} + \text{O}_2 \rightarrow 2\text{MgO} \qquad 4\text{Al} + 3\text{O}_2 \rightarrow 2\text{Al}_2\text{O}_3
\]

\[
\text{P}_4 + 5\text{O}_2 \rightarrow \text{P}_4\text{O}_{10} \qquad \text{S} + \text{O}_2 \rightarrow \text{SO}_2
\]

With \textbf{chlorine}:

\[
2\text{Na} + \text{Cl}_2 \rightarrow 2\text{NaCl} \qquad \text{Mg} + \text{Cl}_2 \rightarrow \text{MgCl}_2 \qquad 2\text{Al} + 3\text{Cl}_2 \rightarrow 2\text{AlCl}_3
\]

\[
\text{Si} + 2\text{Cl}_2 \rightarrow \text{SiCl}_4 \qquad 2\text{P} + 5\text{Cl}_2 \rightarrow 2\text{PCl}_5
\]

With \textbf{water} (only Na and Mg react):

\[
2\text{Na} + 2\text{H}_2\text{O} \rightarrow 2\text{NaOH} + \text{H}_2 \qquad \text{Mg} + 2\text{H}_2\text{O} \rightarrow \text{Mg(OH)}_2 + \text{H}_2
\]

Sodium reacts fast; magnesium reacts only very slowly with cold water.

\subsection*{Oxidation number of the oxides and chlorides}

The \textbf{oxidation number} 氧化数 of the Period 3 element in its oxide or chloride rises across the period, because it equals the number of outer-shell (\textbf{valence shell} 价层) \textbf{electrons} 电子 the atom uses in bonding:

\begin{center}
\adjustbox{max width=\linewidth}{%
\begin{tabular}{llllll}
\toprule
\textbf{Oxides} & \textbf{$\text{Na}_2\text{O}$} & \textbf{$\text{MgO}$} & \textbf{$\text{Al}_2\text{O}_3$} & \textbf{$\text{P}_4\text{O}_{10}$} & \textbf{$\text{SO}_2$ / $\text{SO}_3$} \\
\midrule
oxidation number & $+1$ & $+2$ & $+3$ & $+5$ & $+4$ / $+6$ \\
\bottomrule
\end{tabular}}
\end{center}

The chlorides $\text{NaCl}$, $\text{MgCl}_2$, $\text{AlCl}_3$, $\text{SiCl}_4$, $\text{PCl}_5$ show oxidation numbers $+1$ to $+5$ in the same way.

\subsection*{Oxides with water, and acid–base behaviour}

Across the period the \textbf{oxides} 氧化物 change from basic to acidic:

\begin{center}
\adjustbox{max width=\linewidth}{%
\begin{tabular}{llll}
\toprule
\textbf{Oxide} & \textbf{With water} & \textbf{Acid–base nature} & \textbf{Approximate pH} \\
\midrule
$\text{Na}_2\text{O}$ & $\text{Na}_2\text{O} + \text{H}_2\text{O} \rightarrow 2\text{NaOH}$ & basic & 13–14 \\
$\text{MgO}$ & $\text{MgO} + \text{H}_2\text{O} \rightarrow \text{Mg(OH)}_2$ & basic & 9–10 \\
$\text{Al}_2\text{O}_3$ & insoluble & \textbf{amphoteric} 两性 & 7 \\
$\text{SiO}_2$ & insoluble & weakly acidic & 7 \\
$\text{P}_4\text{O}_{10}$ & $\text{P}_4\text{O}_{10} + 6\text{H}_2\text{O} \rightarrow 4\text{H}_3\text{PO}_4$ & acidic & 1–2 \\
$\text{SO}_3$ & $\text{SO}_3 + \text{H}_2\text{O} \rightarrow \text{H}_2\text{SO}_4$ & strongly acidic & 0–1 \\
\bottomrule
\end{tabular}}
\end{center}

Metal oxides (left) are basic; non-metal oxides (right) are acidic. $\text{Al}_2\text{O}_3$ and its \textbf{hydroxide} 氢氧化物 $\text{Al(OH)}_3$ are amphoteric — they react with \textbf{both} acids and bases:

\[
\text{Al}_2\text{O}_3 + 6\text{HCl} \rightarrow 2\text{AlCl}_3 + 3\text{H}_2\text{O}
\]

\[
\text{Al}_2\text{O}_3 + 2\text{NaOH} + 3\text{H}_2\text{O} \rightarrow 2\text{NaAl(OH)}_4
\]

\subsection*{Chlorides with water}

\begin{itemize}
\item $\text{NaCl}$ and $\text{MgCl}_2$ are ionic. They simply dissolve, giving a near-neutral solution.

\item $\text{SiCl}_4$ and $\text{PCl}_5$ are covalent. They undergo \textbf{hydrolysis} 水解 (react with water) to make acidic solutions and fumes of $\text{HCl}$:

\end{itemize}

\[
\text{SiCl}_4 + 2\text{H}_2\text{O} \rightarrow \text{SiO}_2 + 4\text{HCl} \qquad \text{PCl}_5 + 4\text{H}_2\text{O} \rightarrow \text{H}_3\text{PO}_4 + 5\text{HCl}
\]

\subsection*{Explaining the trends}

These trends follow from the change in bonding and \textbf{electronegativity} 电负性. On the left, the elements are metals with low electronegativity, so their oxides and chlorides are ionic and basic (or neutral). On the right, the elements are non-metals with high electronegativity, so their oxides and chlorides are covalent and acidic. You can use a chloride's or oxide's properties (melting point, conductivity, effect on water) to suggest whether its bonding is ionic or covalent.

\section*{Periodicity of other elements}

The same idea works for any group. If you know the pattern down a group and across a period, you can:

\begin{itemize}
\item \textbf{predict} the properties of an element from its position (for example, a Group 1 element will be a reactive metal forming a $+1$ ion).

\item \textbf{deduce} the likely position and identity of an unknown element from its physical and chemical properties.

\end{itemize}


\end{document}
